% Figura redibujada a partir de "Imagenes/Nivel de Fermi concentracion.PNG" (Tema 1).
% Ocupa lo mismo que el PNG en el documento: 228.2 pt x 153.4 pt. Las
% posiciones se dan en píxeles del PNG original (689 x 463 px).
% Si a 300 K (kT = 0.0259 eV), ionización total:
%   E_C - E_F = kT ln(N_C/N_d),  E_F - E_V = kT ln(N_V/N_a),
%   E_Fi - E_V = E_g/2 - (kT/2) ln(N_C/N_V),  E_g = 1.12 eV.
\documentclass[12pt,border=0pt]{standalone}
\usepackage[spanish,es-tabla]{babel}
\usepackage[utf8]{inputenc}
\usepackage{amsmath}
\usepackage{xcolor}
\usepackage{tikz}
\usepackage{pgfplots}
\usepackage{circuitikz}
\usetikzlibrary{babel, circuits.ee}
\pgfplotsset{compat=1.18}
% Ejes comunes; esquina inferior izquierda en (63, 388) px
\pgfplotsset{ejes/.style={
    at={(20.866pt,-128.51pt)}, anchor=south west, scale only axis,
    width=195.08pt, height=100.36pt,
    xmin=1e12, xmax=1e18, ymin=0, ymax=1.12,
    xtick={1e12,1e13,1e14,1e15,1e16,1e17,1e18}, xminorticks=false,
    tick align=inside, tick style={line width=0.4pt, black}, major tick length=3pt,
    axis line style={line width=0.4pt},
}}

\definecolor{applegreen}{rgb}{0.55, 0.71, 0.0}
\definecolor{airforceblue}{rgb}{0.36, 0.54, 0.66}
\definecolor{amber}{rgb}{1.0, 0.49, 0.0}

\newcommand{\diff}[3][1]{%
	\IfBlankTF{#1}{\frac{\mathrm{d}{#2}}{\mathrm{d}{#3}}}%
	{\frac{\mathrm{d}^{#1} {#2}}{\mathrm{d}{#3}^{#1}}}%
}%
\newcommand{\pdiff}[3][1]{%
	\IfBlankTF{#1}{\frac{\partial {#2}}{\partial {#3}}}%
	{\frac{\partial^{#1} {#2}}{\partial {#3}^{#1}}}%
}%

\begin{document}
\begin{tikzpicture}[x=0.33121pt, y=-0.33121pt, font=\scriptsize]
    \useasboundingbox (0,0) rectangle (689,463);
    \begin{scope}[x=1pt, y=1pt]
    % Eje inferior (N_A) y eje de energías
    \begin{axis}[ejes, xmode=log, clip=false,
        axis x line*=bottom, axis y line*=left,
        ytick={0, 0.5472, 1.12}, yticklabels={$E_\mathrm{V}$, $E_\mathrm{Fi}$, $E_\mathrm{C}$},
        ytick style={draw=none},
        xlabel={$N_\mathrm{A}$ (cm$^{-3}$)}, xlabel style={yshift=2pt},
        declare function={kT = 0.0259;},
    ]
        % Bordes de las bandas
        \addplot[black, ultra thick, domain=12:18, samples=2] ({10^x}, {1.12});
        \addplot[black, ultra thick, domain=12:18, samples=2] ({10^x}, {0});
        \addplot[black, thick, dashed, domain=12:18, samples=2] ({10^x}, {0.5472});
        % Nivel de Fermi: tipo N (azul) y tipo P (rojo)
        \addplot[blue, ultra thick, domain=12:18, samples=40]
            ({10^x}, {1.12 - kT*ln(2.8e19/10^x)});
        \addplot[red, ultra thick, domain=12:18, samples=40]
            ({10^x}, {kT*ln(1.04e19/10^x)});
    \end{axis}
    % Eje superior (N_D)
    \begin{axis}[ejes, xmode=log,
        axis x line*=top, hide y axis,
        xlabel={$N_\mathrm{D}$ (cm$^{-3}$)}, xlabel style={yshift=-2pt},
    ]
    \end{axis}
    \end{scope}
    \node[text=blue] at (435,168) {Tipo N};
    \node[text=red] at (433,302) {Tipo P};
\end{tikzpicture}
\end{document}
